\subsection{Divergence, Upper and Lower Limits, and Sequence Estimates}

\begin{definition} \label{def:sequence-diverges-to-infinity} \hlt{Divergence to $\pm\infty$}\\
For a real sequence $(s_n)$,
\begin{align}
s_n\to+\infty &\iff \forall M\in\R,\ \exists N\in\N:\ n\geq N\Rightarrow s_n\geq M, \nonumber\\
s_n\to-\infty &\iff \forall M\in\R,\ \exists N\in\N:\ n\geq N\Rightarrow s_n\leq M. \nonumber
\end{align}
This extends $\to$ notation for divergent real sequences; Definition \ref{def:real-complex-sequence-convergence} is unchanged.
\end{definition}

\begin{proposition} \label{prop:unbounded-real-sequence-has-infinite-subsequence} \hlt{Unbounded Real Sequences Have Infinite Subsequences}\\
A real sequence is unbounded iff some subsequence tends to $+\infty$ or $-\infty$.
\begin{proof}
Set $n_0=0$. If $(s_n)$ is unbounded above, every tail is unbounded above, so choose $n_k>n_{k-1}$ with $s_{n_k}\geq k$; then $s_{n_k}\to+\infty$. The unbounded-below case similarly gives $s_{n_k}\leq-k$, hence $s_{n_k}\to-\infty$.\\
Conversely, a subsequence tending to $\pm\infty$ is unbounded, so $(s_n)$ is unbounded.\\
\hlt{Structural Commentary}: Unbounded tails supply the thresholds $\pm k$.\\
\hlt{Key Insight}: Large late terms form a divergent subsequence.
\end{proof}
\end{proposition}

\begin{figure}[H]
  \centering
  \begin{tikzpicture}[
  >=Stealth,
  x=0.9cm,
  y=0.9cm,
  every node/.style={font=\small},
  axis/.style={->, gray!70, line width=0.55pt},
  zerodot/.style={circle, fill=orange!85!black, inner sep=1.45pt},
  highdot/.style={circle, fill=blue!70!black, inner sep=1.45pt},
  guide/.style={dashed, line width=0.55pt}
]
  \draw[axis] (0.55,0) -- (7.7,0) node[right] {$n$};
  \draw[axis] (0.80,-0.16) -- (0.80,3.55) node[above] {$s_n$};

  \foreach \x in {1.20,2.55,3.90,5.25,6.60} {
    \node[zerodot] at (\x,0) {};
  }
  \foreach \x/\y in {1.85/0.80,3.20/1.35,4.55/1.90,5.90/2.45,7.25/3.00} {
    \node[highdot] at (\x,\y) {};
  }

  \draw[guide, blue!60!black] (1.85,0.80) -- (7.25,3.00);
  \draw[guide, orange!85!black] (0.95,0) -- (6.85,0);
\end{tikzpicture}

  \caption{For $s_n=[(-1)^n+1]n$, odd terms (orange) stay at $0$, while even terms (blue) diverge to $+\infty$.}
  \label{fig:limsup-extended-subsequences}
\end{figure}

\begin{definition} \label{def:upper-and-lower-limits} \hlt{Upper and Lower Limits}\\
For real sequence $(s_n)$, let $E\defeq\{x\in\overline{\R}:\exists\, n_1<n_2<\cdots,\ s_{n_k}\to x\}$, allowing $x=\pm\infty$.\\
Set $s^*\defeq\sup E$ and $s_*\defeq\inf E$; equivalently, $\limsup_{n\to\infty}s_n=s^*$ and $\liminf_{n\to\infty}s_n=s_*$.
\end{definition}

\begin{theorem} \label{thm:limsup-characterization} \hlt{Characterization of the Upper Limit}\\
Let $(s_n)$ be a real sequence, and let $E,s^*,s_*$ be as in Definition \ref{def:upper-and-lower-limits}. Then:
\begin{enumerate}[label=(\roman*)]
\setlength{\itemsep}{0pt}
\item $s^*\in E$.
\item If $x>s^*$, then $s_n<x$ eventually.
\item These two properties characterize $s^*$ uniquely.
\end{enumerate}
Dually, $s_*$ is uniquely characterized by $s_*\in E$ and $x<s_*$ implies $s_n>x$ eventually.
\begin{proof}
\begin{enumerate}[label=(\roman*)]
\setlength{\itemsep}{0pt}
\item Since $s^*=\sup E$: if $s^*=+\infty$, use Proposition \ref{prop:unbounded-real-sequence-has-infinite-subsequence}; if $s^*\in\R$, use Theorem \ref{thm:subsequential-limits-closed}. If $s^*=-\infty$ and $s_n\nrightarrow-\infty$, infinitely many terms are bounded below; they are either unbounded above or have a bounded subsequence, giving an element of $E$ above $-\infty$, a contradiction.
\item If $s_n\geq x>s^*$ infinitely often, that subsequence is unbounded above or has a finite subsequential limit $y\geq x$; both contradict $s^*=\sup E<x$. Hence $s_n<x$ eventually.
\item If $p<q$ both satisfy (i)--(ii), choose $p<x<q$. Then $s_n<x$ eventually, contradicting $q\in E$.
\end{enumerate}
Apply the same argument to $(-s_n)$ for $s_*$.\\
\hlt{Structural Commentary}: Limsup is the maximal subsequential limit and eventual ceiling.\\
\hlt{Key Insight}: Limsup combines subsequences with eventual upper bounds.
\end{proof}
\end{theorem}

\begin{figure}[H]
  \centering
  \begin{tikzpicture}[
  >=Stealth,
  x=0.9cm,
  y=0.9cm,
  every node/.style={font=\small},
  axis/.style={->, gray!70, line width=0.55pt},
  upper/.style={blue!70!black, line width=0.75pt},
  lower/.style={orange!85!black, line width=0.75pt},
  udot/.style={circle, fill=blue!70!black, inner sep=1.25pt},
  ldot/.style={circle, fill=orange!85!black, inner sep=1.25pt},
  limit/.style={green!45!black, dashed, line width=0.55pt}
]
  \draw[axis] (0.55,0) -- (7.55,0) node[right] {$N$};
  \draw[axis] (0.85,-1.55) -- (0.85,1.65);

  \draw[limit] (1.00,0.62) -- (7.20,0.62);
  \draw[limit] (1.00,-0.62) -- (7.20,-0.62);
  \node[green!45!black, anchor=west] at (7.25,0.62) {$s^*$};
  \node[green!45!black, anchor=west] at (7.25,-0.62) {$s_*$};

  \coordinate (u1) at (1.25,1.34);
  \coordinate (u2) at (2.15,1.04);
  \coordinate (u3) at (3.05,0.86);
  \coordinate (u4) at (3.95,0.75);
  \coordinate (u5) at (4.85,0.68);
  \coordinate (u6) at (5.75,0.64);
  \coordinate (u7) at (6.65,0.625);
  \draw[upper] plot[smooth] coordinates {(u1) (u2) (u3) (u4) (u5) (u6) (u7)};
  \node[blue!70!black, anchor=south] at (2.15,1.12) {$u_N$};

  \coordinate (l1) at (1.25,-1.34);
  \coordinate (l2) at (2.15,-1.03);
  \coordinate (l3) at (3.05,-0.85);
  \coordinate (l4) at (3.95,-0.74);
  \coordinate (l5) at (4.85,-0.68);
  \coordinate (l6) at (5.75,-0.64);
  \coordinate (l7) at (6.65,-0.625);
  \draw[lower] plot[smooth] coordinates {(l1) (l2) (l3) (l4) (l5) (l6) (l7)};
  \node[orange!85!black, anchor=north] at (2.15,-1.11) {$\ell_N$};

  \foreach \p in {u1,u2,u3,u4,u5,u6,u7} {
    \node[udot] at (\p) {};
  }
  \foreach \p in {l1,l2,l3,l4,l5,l6,l7} {
    \node[ldot] at (\p) {};
  }
\end{tikzpicture}

  \caption{Blue tail suprema $u_N$ decrease to $s^*=\limsup s_n$; orange tail infima $\ell_N$ increase to $s_*=\liminf s_n$.}
  \label{fig:limsup-liminf-tail-bounds}
\end{figure}

\begin{proposition} \label{prop:limsup-liminf-tail-suprema} \hlt{Tail Suprema Formula}\\
For a real sequence $(s_n)$, put $T_N\defeq\{s_n:n\geq N\}$ and, in $\overline{\R}$, $u_N\defeq\sup T_N$, $\ell_N\defeq\inf T_N$.\\
Then $u_N\downarrow$, $\ell_N\uparrow$, and $\limsup\limits_{n\to\infty}s_n=\inf_Nu_N$ and $\liminf\limits_{n\to\infty}s_n=\sup_N\ell_N$.
\begin{proof}
Since $T_{N+1}\subseteq T_N$, $u_N\downarrow$ and $\ell_N\uparrow$. Let $s^*=\limsup s_n$ and $u=\inf_Nu_N$.\\
If $x>s^*$, then $s_n<x$ eventually, so $u\leq u_N\leq x$ for some $N$. If $y<s^*$, then $s^*\in E$ gives a subsequence tending to $s^*$; hence every tail meets $(y,\infty)$, so $u_N\geq y$ for all $N$ and $u\geq y$. Thus $u=s^*$.\\
Apply the same formula to $(-s_n)$ and use Corollary \ref{cor:infimum-as-negative-supremum}.\\
\hlt{Structural Commentary}: Tail envelopes converge to $\limsup$ and $\liminf$.\\
\hlt{Key Insight}: Shrinking tails compute the extremal subsequential limits.
\end{proof}
\end{proposition}

\begin{theorem} \label{thm:convergence-iff-limsup-liminf-equal} \hlt{Convergence and Equal Upper/Lower Limits}\\
A real sequence $(s_n)$ converges to a finite number $s\in\R$ if and only if $\limsup_{n\to\infty}s_n=\liminf_{n\to\infty}s_n=s$.
\begin{proof}
\begin{enumerate}[label=\roman*.]
\setlength{\itemsep}{0pt}
\item[$(\Rightarrow)$] If $s_n\to s$, every subsequence converges to $s$, and the tail is bounded, so no subsequence tends to $\pm\infty$. Thus $E=\{s\}$, hence $\limsup s_n=\liminf s_n=s$.
\item[$(\Leftarrow)$] Suppose $\limsup s_n=\liminf s_n=s$. Given $\varepsilon>0$, Theorem \ref{thm:limsup-characterization} gives $s_n<s+\varepsilon$ and $s_n>s-\varepsilon$ eventually. Hence $\abs{s_n-s}<\varepsilon$ eventually, so $s_n\to s$.
\end{enumerate}
\hlt{Structural Commentary}: Equality of the two extremal subsequential limits traps the whole tail.\\
\hlt{Key Insight}: Limsup supplies an eventual ceiling; liminf supplies an eventual floor.
\end{proof}
\end{theorem}

\begin{theorem} \label{thm:limsup-liminf-order-comparison} \hlt{Order Comparison for Upper and Lower Limits}\\
If $s_n\leq t_n$ for all $n\geq N$, where $N$ is fixed, then:
\begin{enumerate}[label=(\roman*)]
\setlength{\itemsep}{0pt}
\item $\limsup\limits_{n\to\infty}s_n\leq \limsup\limits_{n\to\infty}t_n$.
\item $\liminf\limits_{n\to\infty}s_n\leq \liminf\limits_{n\to\infty}t_n$.
\end{enumerate}
\begin{proof}
Discard finitely many terms, so $s_n\leq t_n$ for all $n$.
\begin{enumerate}[label=(\roman*)]
\setlength{\itemsep}{0pt}
\item Put $a=\limsup s_n$ and $b=\limsup t_n$. If $x<a$, then $s_n>x$ infinitely often, hence $t_n>x$ infinitely often; Theorem \ref{thm:limsup-characterization} gives $x\leq b$. Thus $a\leq b$.
\item Apply (i) to $-t_n\leq -s_n$ to get $\limsup(-t_n)\leq\limsup(-s_n)$, equivalently $\liminf s_n\leq\liminf t_n$.
\end{enumerate}
\hlt{Structural Commentary}: Eventual order is inherited by the extremal subsequential limits.\\
\hlt{Key Insight}: Limsup and liminf ignore finitely many initial terms.
\end{proof}
\end{theorem}

\begin{theorem} \label{thm:pinching-theorem-sequences} \hlt{Pinching Theorem for Real Sequences}\\
Let $(a_n)$, $(s_n)$, and $(b_n)$ be real sequences with $a_n\leq s_n\leq b_n$ for all sufficiently large $n$.\\
If $a_n\to s$ and $b_n\to s$, then $s_n\to s$.
\begin{proof}
Let $\varepsilon>0$. Choose $N_0$ s.t. $a_n\leq s_n\leq b_n$ for $n\geq N_0$; $N_1$ s.t. $s-\varepsilon<a_n$ for $n\geq N_1$; and $N_2$ s.t. $b_n<s+\varepsilon$ for $n\geq N_2$. If $N=\max\{N_0,N_1,N_2\}$, then $s-\varepsilon<s_n<s+\varepsilon$ for all $n\geq N$. Hence $s_n\to s$.\\
\hlt{Structural Commentary}: Two convergent bounds squeeze the middle sequence to the same limit.\\
\hlt{Key Insight}: Once both walls are inside the $\varepsilon$-neighborhood, the middle term has no room to escape.
\end{proof}
\end{theorem}

\begin{figure}[H]
  \centering
  \begin{tikzpicture}[
  >=Stealth,
  x=0.9cm,
  y=0.9cm,
  every node/.style={font=\small},
  axis/.style={->, gray!70, line width=0.55pt},
  upper/.style={blue!70!black, line width=0.9pt},
  lower/.style={orange!85!black, line width=0.9pt},
  upperdot/.style={circle, fill=blue!70!black, inner sep=1.35pt},
  lowerdot/.style={circle, fill=orange!85!black, inner sep=1.35pt},
  middledot/.style={circle, draw=green!45!black, fill=white, line width=0.65pt, inner sep=1.45pt}
]
  \draw[axis] (0.35,0) -- (7.55,0) node[right] {$n$};
  \draw[axis] (0.7,-1.45) -- (0.7,1.55);
  \node[left] at (0.65,0) {$s$};

  \coordinate (u1) at (1,1.18);
  \coordinate (u2) at (2,0.72);
  \coordinate (u3) at (3,0.45);
  \coordinate (u4) at (4,0.28);
  \coordinate (u5) at (5,0.18);
  \coordinate (u6) at (6,0.10);
  \coordinate (u7) at (7,0.05);

  \draw[upper]
    plot[smooth] coordinates {(u1) (u2) (u3) (u4) (u5) (u6) (u7)};

  \coordinate (l1) at (1,-1.08);
  \coordinate (l2) at (2,-0.60);
  \coordinate (l3) at (3,-0.34);
  \coordinate (l4) at (4,-0.20);
  \coordinate (l5) at (5,-0.11);
  \coordinate (l6) at (6,-0.06);
  \coordinate (l7) at (7,-0.03);

  \draw[lower]
    plot[smooth] coordinates {(l1) (l2) (l3) (l4) (l5) (l6) (l7)};

  \foreach \p in {u1,u2,u3,u4,u5,u6,u7} {
    \node[upperdot] at (\p) {};
  }
  \foreach \p in {l1,l2,l3,l4,l5,l6,l7} {
    \node[lowerdot] at (\p) {};
  }

  \foreach \x/\ym in {
    1/-0.42,
    2/0.16,
    3/0.02,
    4/0.12,
    5/-0.02,
    6/0.03,
    7/0.01
  } {
    \node[middledot] at (\x,\ym) {};
  }

  \node[blue!70!black, anchor=west] at (1.08,1.20) {$(b_n)$};
  \node[green!45!black, anchor=west] at (1.08,-0.42) {$(s_n)$};
  \node[orange!85!black, anchor=west] at (1.08,-1.2) {$(a_n)$};
\end{tikzpicture}

  \caption{With $a_n\leq s_n\leq b_n$ and $a_n,b_n\to s$, the middle sequence also tends to $s$.}
  \label{fig:pinching-theorem-sequences}
\end{figure}

\begin{theorem} \label{thm:special-sequence-limits} \hlt{Special Sequence Limits}
\begin{enumerate}[label=(\roman*)]
\setlength{\itemsep}{0pt}
\item If $p>0$, then $\lim\limits_{n\to\infty}n^{-p}=0$.
\item If $p>0$, then $\lim\limits_{n\to\infty}\sqrt[n]{p}=1$.
\item $\lim\limits_{n\to\infty}\sqrt[n]{n}=1$.
\item If $p>0$ and $\alpha\in\R$, then $\lim\limits_{n\to\infty}\frac{n^\alpha}{(1+p)^n}=0$.
\item If $\abs{x}<1$, then $\lim\limits_{n\to\infty}x^n=0$.
\end{enumerate}
\begin{proof}
\begin{enumerate}[label=(\roman*)]
\setlength{\itemsep}{0pt}
\item Given $\varepsilon>0$, choose $N>(1/\varepsilon)^{1/p}$ by Theorem \ref{thm:archimedean-and-rational-density}. Then $n\geq N$ gives $n^{-p}<\varepsilon$.
\item If $p>1$, set $x_n\defeq\sqrt[n]{p}-1$. Then $x_n>0$ and $1+nx_n\leq(1+x_n)^n=p$, so $0<x_n\leq(p-1)/n\to0$.\\
By Theorem \ref{thm:pinching-theorem-sequences}, $x_n\to0$. If $p=1$, then $\sqrt[n]{p}=1$ for every $n$.\\
If $0<p<1$, then $1/p>1$, so $\sqrt[n]{1/p}\to1$ by the previous case.\\
Since $\sqrt[n]{p}=1/\sqrt[n]{1/p}$ and limits preserve reciprocals away from $0$, $\sqrt[n]{p}\to1$.
\item Put $x_n\defeq\sqrt[n]{n}-1$. Then $x_n\geq0$, and for $n\geq2$ the binomial theorem gives $n=(1+x_n)^n\geq \frac{n(n-1)}{2}x_n^2$.\\
Thus $0\leq x_n\leq\sqrt{2/(n-1)}$, so Theorem \ref{thm:pinching-theorem-sequences} gives $x_n\to0$.
\item Choose an integer $k>\alpha$ with $k>0$. For $n>2k$, $(1+p)^n\geq \binom{n}{k}p^k=\frac{n(n-1)\cdots(n-k+1)}{k!}p^k>\frac{n^k p^k}{2^k k!}$.\\
Hence $0<\frac{n^\alpha}{(1+p)^n}<\frac{2^kk!}{p^k}n^{\alpha-k}$, and $n^{\alpha-k}\to0$ by (i). Apply Theorem \ref{thm:pinching-theorem-sequences}.
\item If $x\neq0$, write $\abs{x}=(1+p)^{-1}$ with $p>0$ and apply (iv) with $\alpha=0$ to get $\abs{x}^n\to0$. Since $-\abs{x}^n\leq x^n\leq\abs{x}^n$, Theorem \ref{thm:pinching-theorem-sequences} gives $x^n\to0$; $x=0$ is immediate.
\end{enumerate}
\hlt{Structural Commentary}: Each limit is reduced to an elementary comparison sequence.\\
\hlt{Key Insight}: The binomial theorem supplies the right upper bounds for roots and exponential domination.
\end{proof}
\end{theorem}

\begin{method} \label{meth:sequence-proof-patterns}\label{meth:sequence-proof-patterns-source-todo} \hlt{Sequence Proof Patterns}
\begin{enumerate}[label=(\roman*)]
\setlength{\itemsep}{0pt}
\item \hlt{Convergence}: fix $\varepsilon>0$, then find $N$ so that all later terms lie in the relevant $\varepsilon$-ball.
\item \hlt{Uniqueness}: compare two proposed limits with $\varepsilon/2$ and use the triangle inequality.
\item \hlt{Subsequences}: choose $n_k\uparrow$ recursively with distances $<1/k$.
\item \hlt{Cauchy arguments}: estimate the diameter of tails rather than guessing the limit first.
\item \hlt{Monotone sequences}: use $\sup$ or $\inf$ and then prove eventual trapping.
\end{enumerate}
\hlt{Key Insight}: Most sequence proofs are bookkeeping for quantifiers, indices, and one controlling inequality.
\end{method}
